% OWNER: Member 3. DRAFT: every \todo{} must be filled from results/<tag>/ (scripts/analyze.py).
\section{h-DQN: baseline and improvements}\label{sec:hdqn}

\paragraph{Protocol.} Same as the other agents (Section~\ref{sec:setup}): the 15
development games, seeds $\{0,1,2\}$, 5000 actions per run, the 8$\times$8 click
grid, and unmodified official scoring (\todo{arc-agi version, local scorecard}).
Everything is re-initialised at the start of each game; weights, replay memories,
the GRU state and novelty counts persist across levels and resets \emph{within} a
game. Environment reward is $+1$ per level completed; the intrinsic subgoal reward
and the novelty bonus never enter a reported score. The h-DQN code is our own
implementation of \citet{kulkarni2016hdqn} (\texttt{agents/hdqn.py}); no external
h-DQN code or pretrained model is used.

\subsection{Baseline}\label{sec:hdqn-base}
\textbf{Design.} Subgoals are the cells of a 4$\times$4 partition of the 64$\times$64
frame ($|\mathcal G|=16$, 16-pixel cells). Subgoal $g$ is \emph{reached} when at
least one pixel inside cell $g$ changes; it uses only the public frame. In a game
where actions move an avatar this means ``go to cell $g$''; where clicks toggle
objects it means ``act on what is in cell $g$''. The horizon is $H=10$ actions,
$\gamma=0.99$, $\gamma_{\text{m}}=0.95$. Controller and meta-controller use the
shared CNN encoder (the controller gets one extra subgoal plane), and
$\epsilon$ decays from 1 to 0.05 over 3000 actions (controller) and from 1 to 0.1
over 300 segments (meta-controller). Hyperparameters not listed are in
Table~\ref{tab:hparams}.

\textbf{Sanity checks and reference.} The random policy and the flat DQN of
Section~\ref{sec:dqn} (same games, seeds, budgets).

\textbf{Expected limitation.} The meta-controller receives non-zero reward on
very few segments (random play completes 9 levels in 45 runs), so $Q_2$ has almost
nothing to learn from, whereas the controller does get a dense signal \emph{if}
subgoals are reachable. \emph{Hypothesis H0:} the learned meta-controller adds
nothing over choosing subgoals uniformly at random. We test H0 with the
\texttt{random\_meta} ablation, in which the meta-controller is replaced by
uniform sampling and the controller is identical.

\textbf{Result.} \todo{Table~\ref{tab:hdqn-ablation} rows ``h-DQN random meta''
and ``h-DQN''; fraction of subgoals reached (\texttt{goal\_success} in
\texttt{runs.jsonl} diagnostics) over the run; whether H0 holds. If
\texttt{goal\_success} stays near 0 the controller has no usable signal and the
subgoal definition is the bottleneck, which we state as the finding.}

\subsection{Improvement 1: novelty bonus on the meta-controller}\label{sec:hdqn-nov}
\textbf{Limitation $\to$ hypothesis.} \todo{From H0 result:} the meta reward is
almost always 0, so $Q_2$ stays flat and subgoal choice stays near uniform
(measured by the entropy of chosen subgoals, \texttt{goal\_entropy}). We
hypothesise that rewarding subgoals whose segments end in a frame not seen before
in this game gives $Q_2$ a dense, task-agnostic signal that favours subgoals that
change the world, and thereby reaches new states and levels sooner.

\textbf{Intervention.} At the end of each segment the meta reward becomes
$\sum f + \beta/\sqrt{N(h(s_{\text{end}}))}$, where $h$ is a hash of the final frame,
$N$ counts how often it was seen in this game, and $\beta=0.5$
\citep{bellemare2016count}. The controller's intrinsic reward is unchanged. This
differs from flat-DQN novelty because the bonus attaches to \emph{subgoal choices}
rather than to every primitive step; its effect is evaluated separately here.

\textbf{Result and analysis.} \todo{Table~\ref{tab:hdqn-ablation}; change in goal
entropy, number of distinct frames visited, levels gained or lost per game;
possible failure mode: frames that change every step (timers, animations) make
every segment look novel.}

\subsection{Improvement 2: recurrent (GRU) meta-controller}\label{sec:hdqn-gru}
\textbf{Limitation $\to$ hypothesis.} The meta-controller sees one frame, but the
rules are hidden and the right next subgoal depends on what has already been
tried in this game. \todo{Cite a failure replay in which the agent repeats a
fruitless subgoal.} We hypothesise that a memory over past (subgoal, outcome)
pairs lets $Q_2$ stop repeating subgoals that did nothing.

\textbf{Intervention.} $Q_2$ becomes a GRU \citep{cho2014gru} over the sequence of
meta decisions: at each decision its input is the encoded frame, the one-hot
previous subgoal and the previous segment's outcome (reached flag, length$/H$),
hidden size 128. The hidden state persists across levels and resets within a game
and is reset between games. Training follows the recurrent-DQN recipe
\citep{hausknecht2015drqn}: windows of 9 consecutive stored segments, unrolled from a
zero state, with the first 2 steps as burn-in (no loss), as in
\citet{kapturowski2019r2d2}; no recurrent state is stored. The non-recurrent
variant is trained on the same windows with loss on all positions, so the
\emph{only} difference is the GRU. Compute is higher (a sequence of
frames per update).

\textbf{Result and analysis.} \todo{Table~\ref{tab:hdqn-ablation}; learning
curves; whether memory helps, only in some games, or not at all given $\approx$230
meta updates per run.}

\subsection{Ablation and attribution}
\begin{table}[t]
\centering\footnotesize\setlength{\tabcolsep}{3pt}
\caption{h-DQN versions on the 15 development games, 3 seeds, 5000 actions/run.
Official local score in \% (mean $\pm$ sd over seeds), total levels over 45 runs,
and mean wall-clock per run. \todo{fill from \texttt{scripts/analyze.py}.}}
\label{tab:hdqn-ablation}
\begin{tabular}{lcccccc}
\toprule
Version & Meta & Novelty & GRU & Score (\%) & Levels & Wall (s) \\
\midrule
Random             & --         & --         & --         & 0.0104 $\pm$ 0.011 & 9 & \todo{} \\
Flat DQN (Sec.~\ref{sec:dqn}) & -- & -- & -- & \todo{} & \todo{} & \todo{} \\
h-DQN, random meta & random     & $\times$   & $\times$   & \todo{} & \todo{} & \todo{} \\
h-DQN (baseline)   & learned    & $\times$   & $\times$   & \todo{} & \todo{} & \todo{} \\
+ Imp.\ 1 (novelty)& learned    & $\surd$    & $\times$   & \todo{} & \todo{} & \todo{} \\
+ Imp.\ 2 (GRU)    & learned    & $\times$   & $\surd$    & \todo{} & \todo{} & \todo{} \\
+ Imp.\ 1 + 2      & learned    & $\surd$    & $\surd$    & \todo{} & \todo{} & \todo{} \\
\bottomrule
\end{tabular}
\end{table}
The design isolates the hierarchy (random meta vs.\ learned meta), each
improvement alone, and their interaction under identical games, seeds and
budgets. \todo{Learning curves (best levels vs.\ actions, one line per seed);
paired per-game differences; what the hierarchy buys over flat DQN; which effects
are distinguishable from seed noise with three seeds; held-out results run by
Member~1.}

\subsection{Failure cases and limitations}
\todo{Replays of one success and one failure (\texttt{record: true}); games
where no version progressed; honest statement that at 5000 actions the agent has
had roughly 1100 controller and 230 meta updates, so conclusions concern the
early-learning regime. Known simplifications: one decaying $\epsilon$ per level
instead of per-subgoal annealing; time-outs are not terminal for the controller;
GAME\_OVER resets are ordinary transitions; GRU windows start from a zero state;
subgoals are a fixed partition of the screen, not objects.}
